\documentclass[10pt,a4paper]{amsart}
\usepackage{amsmath,amssymb,amsthm,graphicx}
\usepackage{verbatim}
\newenvironment{reference}{\comment}{\endcomment}
%\newenvironment{reference}{}{}
\newenvironment{slogan}{\comment}{\endcomment}
\newenvironment{history}{\comment}{\endcomment}

% For commutative diagrams we use Xy-pic
\usepackage[all]{xy}

% We use 2cell for 2-commutative diagrams.
\xyoption{2cell}
% Theorem environments.
%
\theoremstyle{plain}
\newtheorem{theorem}[subsection]{Theorem}
\newtheorem{proposition}[subsection]{Proposition}
\newtheorem{lemma}[subsection]{Lemma}

\theoremstyle{definition}
\newtheorem{definition}[subsection]{Definition}
\newtheorem{example}[subsection]{Example}
\newtheorem{exercise}[subsection]{Exercise}
\newtheorem{situation}[subsection]{Situation}

\theoremstyle{remark}
\newtheorem{remark}[subsection]{Remark}
\newtheorem{remarks}[subsection]{Remarks}

\numberwithin{equation}{subsection}

% Macros
%
\def\lim{\mathop{\mathrm{lim}}\nolimits}
\def\colim{\mathop{\mathrm{colim}}\nolimits}
\def\Spec{\mathop{\mathrm{Spec}}}
\def\Hom{\mathop{\mathrm{Hom}}\nolimits}
\def\Ext{\mathop{\mathrm{Ext}}\nolimits}
\def\SheafHom{\mathop{\mathcal{H}\!\mathit{om}}\nolimits}
\def\SheafExt{\mathop{\mathcal{E}\!\mathit{xt}}\nolimits}
\def\Sch{\mathit{Sch}}
\def\Mor{\mathop{\mathrm{Mor}}\nolimits}
\def\Ob{\mathop{\mathrm{Ob}}\nolimits}
\def\Sh{\mathop{\mathit{Sh}}\nolimits}
\def\NL{\mathop{N\!L}\nolimits}
\def\CH{\mathop{\mathrm{CH}}\nolimits}
\def\proetale{{pro\text{-}\acute{e}tale}}
\def\etale{{\acute{e}tale}}
\def\QCoh{\mathit{QCoh}}
\def\Ker{\mathop{\mathrm{Ker}}}
\def\Im{\mathop{\mathrm{Im}}}
\def\Coker{\mathop{\mathrm{Coker}}}
\def\Coim{\mathop{\mathrm{Coim}}}

% Boxtimes
%
\DeclareMathSymbol{\boxtimes}{\mathbin}{AMSa}{"02}

%
% Macros for moduli stacks/spaces
%
\def\QCohstack{\mathcal{QC}\!\mathit{oh}}
\def\Cohstack{\mathcal{C}\!\mathit{oh}}
\def\Spacesstack{\mathcal{S}\!\mathit{paces}}
\def\Quotfunctor{\mathrm{Quot}}
\def\Hilbfunctor{\mathrm{Hilb}}
\def\Curvesstack{\mathcal{C}\!\mathit{urves}}
\def\Polarizedstack{\mathcal{P}\!\mathit{olarized}}
\def\Complexesstack{\mathcal{C}\!\mathit{omplexes}}
% \Pic is the operator that assigns to X its picard group, usage \Pic(X)
% \Picardstack_{X/B} denotes the Picard stack of X over B
% \Picardfunctor_{X/B} denotes the Picard functor of X over B
\def\Pic{\mathop{\mathrm{Pic}}\nolimits}
\def\Picardstack{\mathcal{P}\!\mathit{ic}}
\def\Picardfunctor{\mathrm{Pic}}
\def\Deformationcategory{\mathcal{D}\!\mathit{ef}}
\usepackage{cleveref}
\crefname{originallemma}{original lemma}{original lemmas}
\Crefname{originallemma}{Original lemma}{Original lemmas}
\makeatletter
\newlabel{cohomology-lemma-section-RHom-over-U}{{08DK}{source}{}{sourcepointer-1}{}}
\newlabel{cohomology-lemma-section-RHom-over-U@cref}{{[originallemma][0][]08DK}{[1][0][]source}{}{}{}}
\makeatother
\crefname{originallemma}{original lemma}{original lemmas}
\Crefname{originallemma}{Original lemma}{Original lemmas}
\makeatletter
\newlabel{derived-lemma-ext-1}{{06XV}{source}{}{sourcepointer-2}{}}
\newlabel{derived-lemma-ext-1@cref}{{[originallemma][0][]06XV}{[1][0][]source}{}{}{}}
\makeatother
\title{The conormal sheaf of a closed subscheme computes the first Ext sheaf of its structure sheaf}
\author{The Stacks Project Authors}
\date{stacks-project:duality.tex:local-snapshot-20261012; GNU FDL 1.2 or later}
\begin{document}
\maketitle
\noindent\textit{Editorial note.} All delivered mathematical text is the author's own, copied from the pinned Stacks Project chapter file duality.tex; no statement or proof step was rewritten, reordered or removed.\par\medskip
\noindent\textit{Editorial note.} Preamble provenance: the pinned chapter file opens with an input of the shared preamble and carries no preamble text of its own, so the project's single Stacks-wide preamble staged at sources/stacks-duality/preamble.tex is excerpted verbatim under the recipe fixed by LANE\_BRIEF 22.1, as three spans, lines 8-13 (verbatim package and the reference/slogan/history environments), lines 14-18 (Xy-pic setup) and lines 154-217 (the native theorem declarations and mathematical macros). Nothing was added, removed or redefined; the wrapper class and the amsmath/amssymb/amsthm/graphicx base packages are supplied by the assembler.\par\medskip
\noindent\textit{Editorial note.} The counted statement and proof contain no cross reference, no citation and no included file, so the delivered body is self contained.\par\medskip
\noindent\textit{Editorial note.} Canonical location of the counted proposition: Stacks tag 0BQW (sources/stacks-duality/pinned-original-source.tex line 3620).\par\medskip
\noindent\textit{Allowed standard prerequisites (editorial).}\begin{itemize}
\item The second order thickening sequence. Same chapter, equation labelled equation-second-order-thickening, used at proof lines 3660-3666. Supplies the extension whose pushouts define the map.
\item Sections of the derived Hom over an open and the interpretation of Ext\textasciicircum{}1 by extensions. Cohomology, Lemma labelled cohomology-lemma-section-RHom-over-U and Derived Categories, Lemma labelled derived-lemma-ext-1, cited at proof lines 3670-3674. Identify pushouts of the thickening with elements of the Ext sheaf.
\end{itemize}\medskip
\ifdefined\hypertarget\hypertarget{sourcepointer-1}{}\fi
\noindent\textit{Source reference.} Original lemma 08DK: sources/stacks-chow/tags.txt, source lines 10640--10640. Cross-chapter standard prerequisite located in the project tag list staged in the local snapshot; invoked at proof line 3644 of the counted proof. It is retained only as a navigation pointer; every load-bearing step is inside the delivered span.\par\medskip
\ifdefined\hypertarget\hypertarget{sourcepointer-2}{}\fi
\noindent\textit{Source reference.} Original lemma 06XV: sources/stacks-chow/tags.txt, source lines 8864--8864. Cross-chapter standard prerequisite located in the project tag list staged in the local snapshot; invoked at proof line 3645 of the counted proof. It is retained only as a navigation pointer; every load-bearing step is inside the delivered span.\par\medskip
\begin{lemma}
\label{lemma-ext}
Let $Y$ be a ringed space. Let $\mathcal{I} \subset \mathcal{O}_Y$
be a sheaf of ideals. Set $\mathcal{O}_X = \mathcal{O}_Y/\mathcal{I}$ and
$\mathcal{N} =
\SheafHom_{\mathcal{O}_Y}(\mathcal{I}/\mathcal{I}^2, \mathcal{O}_X)$.
There is a canonical isomorphism
$c : \mathcal{N} \to
\SheafExt^1_{\mathcal{O}_Y}(\mathcal{O}_X, \mathcal{O}_X)
$.
\end{lemma}

\begin{proof}
Consider the canonical short exact sequence
\begin{equation}
\label{equation-second-order-thickening}
0 \to \mathcal{I}/\mathcal{I}^2 \to \mathcal{O}_Y/\mathcal{I}^2 \to
\mathcal{O}_X \to 0
\end{equation}
Let $U \subset X$ be open and let $s \in \mathcal{N}(U)$. Then we can
pushout (\ref{equation-second-order-thickening}) via $s$ to
get an extension $E_s$ of $\mathcal{O}_X|_U$ by $\mathcal{O}_X|_U$.
This in turn defines a section $c(s)$ of
$\SheafExt^1_{\mathcal{O}_Y}(\mathcal{O}_X, \mathcal{O}_X)$
over $U$.
See Cohomology, Lemma \ref{cohomology-lemma-section-RHom-over-U}
and Derived Categories, Lemma \ref{derived-lemma-ext-1}.
Conversely, given an extension
$$
0 \to \mathcal{O}_X|_U \to \mathcal{E} \to \mathcal{O}_X|_U \to 0
$$
of $\mathcal{O}_U$-modules, we can find an open covering
$U = \bigcup U_i$ and sections $e_i \in \mathcal{E}(U_i)$
mapping to $1 \in \mathcal{O}_X(U_i)$. Then $e_i$ defines a map
$\mathcal{O}_Y|_{U_i} \to \mathcal{E}|_{U_i}$ whose kernel
contains $\mathcal{I}^2$. In this way we see that
$\mathcal{E}|_{U_i}$ comes from a pushout as above.
This shows that $c$ is surjective. We omit the proof
of injectivity.
\end{proof}

\end{document}
